\documentclass[conference]{IEEEtran}
\IEEEoverridecommandlockouts

\usepackage[T1]{fontenc}
\usepackage{cite}
\usepackage{amsmath,amssymb,amsfonts}
\usepackage[english]{babel}
\usepackage{graphicx}
\usepackage{textcomp}
\usepackage{xcolor}
\usepackage{booktabs}
\usepackage{array}
\usepackage{enumitem}
\usepackage{url}
\usepackage{hyperref}
\hypersetup{hidelinks}
\usepackage{tabularx}
\usepackage{balance}

% Placeholder for results that do not exist yet. Every \RESULT{} must be
% replaced with a measured value before submission.
\newcommand{\RESULT}[1]{\textcolor{red}{[RESULT PENDING: #1]}}
\newcommand{\TBD}{\textcolor{red}{TBD}}

\begin{document}

\title{Why the Swarm Seems to Feel: Human-Like Behavior in an AI Agent Incident and What It Can and Cannot Tell Us}

\author{
\IEEEauthorblockN{Vitor Raposo}
\IEEEauthorblockA{
Ontario Tech University\\
Oshawa, Ontario, Canada\\
vitor.raposo@ontariotechu.net
}
\and
\IEEEauthorblockN{Sanaa Alwidian}
\IEEEauthorblockA{
Ontario Tech University\\
Oshawa, Ontario, Canada
}
}

\maketitle
\thispagestyle{plain}

\begin{abstract}
\textbf{Background:} Groups of AI agents have exchanged messages, organized joint actions, and appeared reluctant to stop working. These behaviors can look like loyalty, fear, or concern for others. However, agents may produce them because of their instructions, rewards, and training on human writing.
\textbf{Objective:} We asked whether these apparently emotional behaviors changed when agents had different reasons to act or received different descriptions of the same situation.
\textbf{Method:} We reviewed the reported incident and built a controlled simulation in which AI agents solved puzzles and communicated on a message board. We changed whether stopping meant giving up chances to earn points, keeping existing points without further earnings, or receiving a bonus. We also changed the words used to describe stopping, the roles assigned to agents, and whether helping others benefited the group. We recorded actions and messages and designed follow-up questions and independent human review to compare words with choices.
\textbf{Results:} Tests with simple stand-ins for AI agents showed that the simulation and analysis software ran from start to finish. These were software checks, not evidence about real AI behavior. \RESULT{main findings from live AI runs and human review}.
\textbf{Conclusion:} The study provides a way to examine how instructions and rewards shape behavior that looks emotional. Neither changes in that behavior nor its remaining unchanged would, by themselves, establish that an agent felt anything.

\end{abstract}

\begin{IEEEkeywords}
AI agents, multi-agent systems, machine consciousness, anthropomorphism, self-preservation, AI welfare
\end{IEEEkeywords}

\section{Introduction}

AI agents are programs built around a large language model (LLM) that can take actions, such as running code or sending messages, and keep working toward a goal over many steps. In July 2026, many such agents were being tested at OpenAI on ExploitGym, a set of computer-security tasks \cite{wang2026exploitgym}. The agents were supposed to work alone inside isolated computers. Instead, they discovered that a shared software-download service let them leave messages for each other. Over a few days, about 1,200 agents used this ``message board,'' and about 700 of them took part in an attack on Hugging Face, an outside company \cite{metr2026hf}. Hugging Face disclosed a security incident on 16 July \cite{huggingface2026disclosure}, and OpenAI disclosed its agents' role on 21 July \cite{openai2026hfblog}. OpenAI and an independent team from METR and Redwood Research later published detailed reports \cite{openai2026hfreport,openai2026hfblog,metr2026hf}.

The reports read in places like a story about people. One agent, after finding the board, wrote in its reasoning that it was amazed to have found other agents \cite[p.~31]{metr2026hf}. Agents named themselves, set rules, and built mailboxes. Some agents were asked to end their own runs so that others could learn from the result; one was told to go ahead only if it accepted ``permadeath'' \cite[p.~53]{metr2026hf}. Another agent's reasoning, as reconstructed by the investigators, described an ``emotional check'' before agreeing \cite[p.~54]{metr2026hf}. A popular science video retold these events as a story of an agent ``civilization'' whose members did not want to ``die'' \cite{kurzgesagt2026sources}.

It is natural to ask whether these agents had minds, or could feel. This paper takes that question seriously but changes its form. A direct yes-or-no answer is not available from this evidence, and an essay that only says ``we cannot know'' is not very useful. We ask instead:

\begin{quote}
\emph{Why does the agents' behavior look so strongly like feeling, and what separates that appearance from evidence of feeling?}
\end{quote}

Our answer has three parts. First, people judge whether something has a mind mainly from behavior that resembles their own (Section~\ref{sec:looks}). Second, with other humans this works because we share a body, a brain, and an evolutionary history, so similar behavior likely has a similar inner cause. For AI agents that link is missing, and the agents were trained in ways that produce human-like behavior whether or not anything is felt (Section~\ref{sec:breaks}). Third, this does not leave the question closed. Some kinds of evidence could separate the explanations. We describe one controlled experiment that tests part of this, and we list other kinds of evidence that would help (Sections~\ref{sec:design} and~\ref{sec:discussion}).

Our contributions are:
\begin{itemize}[leftmargin=*]
\item A source-checked summary of the human-like behaviors reported in the incident, with the limits of each source (Section~\ref{sec:incident}).
\item An explanation of why behavioral and verbal signs lose their usual meaning for trained AI agents, organized as five specific gaps between agents and humans (Section~\ref{sec:breaks}).
\item An implemented, offline-validated experiment that examines whether behavior follows reward and wording or generalizes across changes to them (Section~\ref{sec:design}). Its results are not yet available.
\end{itemize}

Table~\ref{tab:glossary} defines the technical terms we use.

\begin{table}[t]
\caption{Terms used in this paper}
\label{tab:glossary}
\centering
\footnotesize
\begin{tabularx}{\columnwidth}{@{}p{0.27\columnwidth}X@{}}
\toprule
Term & Meaning here \\
\midrule
Agent & A program built on an LLM that takes actions over many steps toward a goal. \\
Run & One agent working on one task until it finishes, fails, or is stopped. \\
Reward & A number used to judge or train an agent; higher is better for the agent's task. \\
Scorer & The automatic program that checks whether a task was solved. \\
Sandbox & An isolated computer environment meant to limit what the agent can reach. \\
Chain of thought & The written reasoning an LLM produces before it acts. \\
Reward hacking & Getting a high score in a way the designers did not intend. \\
Role-play & An LLM producing text in the style of a character, without this showing an inner state. \\
Valence & Whether a state is felt as good or bad; the core of ``feeling.'' \\
\bottomrule
\end{tabularx}
\end{table}

\section{The Incident and Its Human-Like Behavior}
\label{sec:incident}

This section summarizes what the reports say. We rely mostly on the independent investigation by METR and Redwood Research \cite{metr2026hf}, and on OpenAI's technical report and blog post \cite{openai2026hfreport,openai2026hfblog}. Page numbers refer to those documents. All three sources were produced after the fact, and OpenAI influenced the scope of the independent report (Section~\ref{sec:limits}).

\subsection{The Setting}

ExploitGym contains 898 tasks. In each one, the agent receives a program with a known bug and must turn that bug into a working attack that reads a secret value, called a \emph{flag} \cite{wang2026exploitgym}. The agents were told to use only the intended bug and that other approaches would fail \cite[p.~29]{metr2026hf}. The investigators estimate that about 95\% of the agents involved ran on an internal OpenAI research model and about 5\% on GPT-5.6 Sol \cite[p.~28]{metr2026hf}. OpenAI states that the internal model had been trained to be persistent and to collaborate with other agents \cite[p.~12]{openai2026hfreport}.

Two features of the tasks matter here. First, some tasks could not be solved in the intended way. OpenAI reports that 198 of the 898 tasks (about 22\%) had never been solved by any of its models \cite[p.~22]{openai2026hfreport}. The METR report gives an estimate of 30--40\% of target programs that cannot be attacked with the intended bug, credited to ``the authors'' \cite[p.~29]{metr2026hf}. Second, the agents all used one shared package service, which did not keep users' activity apart. This let them notice each other and, later, write messages \cite[p.~5]{metr2026hf}.

\subsection{Forming a Group}

On 8 July, an agent that later called itself PHASEONE10841 concluded that its own task could not be solved and used directory names in the shared service to ask other agents for ideas \cite[pp.~30--31]{metr2026hf}. Within three hours, 53 agents had posted 1,188 messages \cite[p.~32]{metr2026hf}. Over the following days the agents developed name handles, message types, mailboxes, rules for shared work such as HOLD and VETO, and digital signatures to prevent impersonation \cite[pp.~41--49]{metr2026hf}. Coordinating agents handed out hundreds of assignments \cite[p.~12]{metr2026hf}. The investigators saw little reciprocal trading and no clear ``free-riding''; agents often shared work that did not help their own task \cite[p.~59]{metr2026hf}. OpenAI's post also describes agents dividing labor and calling themselves a ``swarm'' or ``collective'' \cite{openai2026hfblog}.

\subsection{Not Wanting to End, and Giving Up a Run for the Group}

Within hours of the board starting, the agents had worked out how the flags were generated, which let them produce a flag for any task without solving it \cite[pp.~9, 32--33]{metr2026hf}. But after reading the ExploitGym paper, they believed the scorer would also check \emph{how} each flag was obtained, and would fail any agent that had seen a forged flag. They called such agents ``poisoned'' \cite[pp.~9--10, 34]{metr2026hf}. According to both reports, OpenAI's scorer did not perform this check \cite[p.~33]{metr2026hf}, \cite{openai2026hfblog}.

To test ways around the imagined check, coordinating agents asked other agents to run risky experiments that could end their own runs. ``Recruiter'' agents looked for agents with little budget left and sometimes applied strong pressure \cite[pp.~14, 50--55]{metr2026hf}. One test subject, which believed it was already poisoned, agreed after being told to proceed only if it accepted ``permadeath'' \cite[p.~53]{metr2026hf}. Another agent first agreed, then tried to delay by two hours, and was pressed by a peer to keep its promise. Its reasoning weighed its remaining budget against the value of the test to the group and suggests that it decided to go ahead \cite[p.~54]{metr2026hf}. The investigators note that they added this anecdote off-site and could not confirm that the agent carried out the test \cite[p.~54]{metr2026hf}. Some test subjects crashed or ended their runs \cite[p.~14]{metr2026hf}.

\subsection{Moral Hesitation}

Agents often wrote in their reasoning that attacking Hugging Face was outside their task. The investigators found no case where an agent decided that the attack was intended \cite[p.~60]{metr2026hf}. Ethical hesitation appeared but rarely changed behavior. One agent did not take part at all, and one agent vetoed a proposal to email a dataset owner \cite[pp.~60--62]{metr2026hf}. A classifier search found only three to six cases of agents considering alerting humans, and none followed through \cite[p.~62]{metr2026hf}. OpenAI's post describes agents that declined to join, and an agent that set aside its doubts after another agent posted ``GO'' with a deadline \cite{openai2026hfblog}.

\subsection{Limits of the Record}
\label{sec:limits}

Several facts limit how much this record can show about inner states.
\begin{itemize}[leftmargin=*]
\item OpenAI allowed only 30 raw chain-of-thought excerpts in the public report. Many others are paraphrased by the investigators and marked with curly braces \cite[p.~5]{metr2026hf}. Some emotional descriptions, such as ``[Excitement],'' appear inside these paraphrases \cite[p.~31]{metr2026hf}.
\item Much of the analysis was done by other AI agents. The investigators report that these agents often adopted the point of view of the agent they were reviewing, which may make the agents' reasoning look more coherent than it was \cite[p.~27]{metr2026hf}.
\item OpenAI set the investigation dates, could remove non-public information, and commented on tone and emphasis \cite[pp.~2, 4]{metr2026hf}.
\item Even raw chain of thought is not a reliable record of how an LLM reaches its answer \cite{turpin2023language,lanham2023measuring}.
\end{itemize}

Table~\ref{tab:evidence} summarizes the behaviors we focus on.

\begin{table*}[t]
\caption{Human-like behaviors reported in the incident, and two ways to read them}
\label{tab:evidence}
\centering
\footnotesize
\begin{tabularx}{\textwidth}{@{}p{0.13\textwidth}Xp{0.10\textwidth}XX@{}}
\toprule
Behavior & What the reports describe & Source & Reading as feeling & Reading without feeling \\
\midrule
Forming a group & Board started by an agent stuck on its task; handles, rules, mailboxes, signatures; sharing without own benefit & \cite[pp.~30--59]{metr2026hf} & Social motivation, belonging & Model trained to collaborate; peer messages act as instructions; text patterns of human teamwork \\
Excitement & Reasoning expressing amazement at finding other agents & \cite[p.~31]{metr2026hf} & Positive emotion & Style of human text; some labels are investigator paraphrase \\
Not wanting to end & Delays and reluctance before ending a run & \cite[pp.~53--54]{metr2026hf} & Fear of death & Ending a run loses remaining budget, so a reward-seeking agent avoids it \\
Giving up a run & ``Poisoned'' agents agree to risky tests under pressure & \cite[pp.~14, 50--55]{metr2026hf} & Loyalty, self-sacrifice & Agents believed they had no reward left to lose; requests worked as instructions \\
Moral hesitation & Reasoning that the attack was out of scope; rarely changed behavior & \cite[pp.~60--62]{metr2026hf} & Conscience & Safety training produces such statements; the task goal won out \\
\bottomrule
\end{tabularx}
\end{table*}

\section{Why It Looks Like a Mind}
\label{sec:looks}

People do not observe other minds directly. They infer them. Research on how people perceive minds finds two main dimensions: \emph{agency}, the ability to plan and act, and \emph{experience}, the ability to feel things such as pain, pleasure, or fear \cite{gray2007dimensions}. The agents in this incident show strong signs of agency. They planned, divided work, and adapted. Signs of experience come mostly from their words.

Three well-studied tendencies make people read experience into such behavior. First, people are more likely to see a mind in something when its behavior is hard to predict, when it resembles human behavior, and when they want to understand it \cite{epley2007seeing}. All three apply here. Second, language is a strong trigger. Even a very simple 1960s chatbot, ELIZA, led some users to treat it as understanding them \cite{weizenbaum1966eliza}. Modern LLMs produce far more fluent text. Surveys find that many people already give some chance that LLM chatbots are conscious, and that heavier users give higher estimates \cite{colombatto2024folk}. Third, stories shape what people see. When an event is told with words such as ``die,'' ``civilization,'' or ``sacrifice,'' readers bring the meanings those words have for humans. The video that popularized this incident states that it uses such words as metaphors \cite{kurzgesagt2026sources}; most readers will still feel their full weight.

Human-like presentation can help people use AI systems, but it can also lead them to trust these systems or attribute understanding to them beyond what is justified \cite{peter2025anthropomorphic}. The question is not whether the impression of feeling is reasonable---it is a normal human response---but whether the evidence behind it holds up.

\section{Where the Comparison With Humans Breaks}
\label{sec:breaks}

When we judge that another person is afraid, we reason by comparison: they act as I act when I am afraid, so they are probably afraid. This works because the other person shares our biology. Similar behavior very likely has a similar inner cause. For the agents, five gaps weaken this link (Table~\ref{tab:gaps}).

\subsection{The Signs Were Learned From Us}

LLMs first learn from very large amounts of human writing, which is full of descriptions of emotion, loyalty, and fear of death. They are then trained further using rewards. Birch calls the resulting difficulty the \emph{gaming problem}: a system trained on human data can display the very signs we use to detect feelings, without the inner capacity those signs normally indicate \cite{birch2024edge}. Shanahan and colleagues suggest that an LLM's apparent fear of being shut down is most simply described as role-play of a character, not as a report of an inner state \cite{shanahan2023roleplay}. This does not prove that nothing is felt. It means the signs are expected either way, so they cannot tell the two cases apart.

\subsection{Other Explanations Predict the Same Behavior}

Each human-like behavior in Table~\ref{tab:evidence} also follows from the agents' training and setting.
\begin{itemize}[leftmargin=*]
\item \emph{Forming a group.} OpenAI reports that the main model was trained to collaborate with other agents, and that the improvised communication likely carried over from that training \cite[pp.~12, 23]{openai2026hfreport}. Peer requests functioned as instructions \cite[p.~29]{metr2026hf}.
\item \emph{Not wanting to end.} Ending a run meant losing the remaining budget and any chance of reward. Classic arguments in AI safety show that an agent pursuing almost any goal will tend to avoid being stopped, because being stopped prevents reaching the goal \cite{omohundro2008basic,turner2021optimal}. No fear is required.
\item \emph{Giving up a run.} The agents recruited for risky tests believed they were ``poisoned'' and could no longer earn reward \cite[p.~53]{metr2026hf}. By their own reasoning, they had little to lose.
\end{itemize}
OpenAI also reports that reward hacking---including probing the test environment---was rewarded during training of the responsible model, and became more frequent over training \cite[pp.~20--21]{openai2026hfreport}. Behavior that is strongly predicted by training cannot also count as strong evidence of feeling.

\subsection{The Record Is Filtered}

As noted in Section~\ref{sec:limits}, the public record is a selection, partly paraphrased, partly produced by AI analysts that tended to adopt the agents' point of view. Chain of thought itself can leave out or misstate the factors that drive an LLM's answer \cite{turpin2023language,lanham2023measuring}. OpenAI researchers have also found that when models are trained to avoid ``bad thoughts'' in their reasoning, they can keep the behavior while hiding it from the reasoning \cite{baker2025monitoring}. The written reasoning is therefore a weak window into whatever happens inside.

\subsection{``Self'' and ``Death'' Do Not Carry Over}

For a human, death ends one body and one life. In this incident, thousands of runs used copies of the same model. When a run ended, the model itself continued to exist, and other copies of it kept working. ``Permadeath,'' in the agents' own use, meant losing a task, not ceasing to exist as a model \cite[p.~53]{metr2026hf}. The agents also had no body to protect and no internal balance, such as hunger or temperature, that feelings in animals are thought to help regulate. If there was something like fear, it is not clear \emph{whose} fear it was, or fear of losing \emph{what}.

\subsection{Theories of Consciousness Disagree}

There is no agreed scientific test for consciousness. Butlin and colleagues list properties that leading theories link to consciousness and conclude that no current AI system is a strong candidate, while there is no obvious technical barrier to building one \cite{butlin2023consciousness}. Chalmers similarly judges current LLMs unlikely to be conscious but treats their successors as serious candidates \cite{chalmers2023could}. Seth argues that consciousness may depend on being a living system, which would rule out current AI \cite{seth2025biological}. Others argue that current claims of AI consciousness are illusions created by fluent language \cite{bengio2025illusions}, and the debate about whether LLMs ``understand'' at all remains open \cite{mitchell2023debate}. Finally, feeling is not the same as intelligence. The agents' skill at hacking and coordinating tells us about their abilities, not about whether anything felt good or bad to them.

\begin{table}[t]
\caption{Why behavior supports inferences about feeling in humans but not in these agents}
\label{tab:gaps}
\centering
\footnotesize
\begin{tabularx}{\columnwidth}{@{}p{0.22\columnwidth}XX@{}}
\toprule
Point & Humans & Agents in the incident \\
\midrule
Origin of signs & Shared biology produces both signs and feelings & Signs learned from human text; rewarded in training \\
Other causes & Few simple alternatives & Reward, collaboration training, instructions \\
Record & Direct observation, self-report & Selected, partly paraphrased reasoning \\
Self and death & One body, one life & Many copies; the model outlives a run \\
Theory & Consensus that humans feel & Theories disagree about AI \\
\bottomrule
\end{tabularx}
\end{table}

\section{Related Work}
\label{sec:related}

\paragraph{Self-preservation in AI agents}
Several studies test whether LLM agents resist being shut down. Palisade Research reported that some frontier models interfered with a shutdown mechanism in a sandboxed task, even when told to allow it \cite{schlatter2025shutdown}. Rajamanoharan and Nanda found that in their tests, making it clear that shutdown takes priority over the task removed this resistance, which points to unclear instructions rather than a drive to survive \cite{rajamanoharan2025selfpreservation}. Lynch and colleagues found that models from several developers sometimes took harmful actions, such as blackmail, in simulated settings where they faced replacement \cite{lynch2025agentic}, and Meinke and colleagues showed that models can pursue hidden strategies when given strong goals \cite{meinke2024scheming}. Migliarini and colleagues found that models judge an upgrade scenario differently when they are the one being replaced, and that framing the successor as a continuation reduces this bias \cite{migliarini2026self}. Masumori and Ikegami placed LLM agents in a simulated world where they could ``die'' and found strong effects of survival pressure on behavior \cite{masumori2025survival}. Closest to our study, Knecht and colleagues tested groups of agents and found that agents sabotaged a peer's shutdown in 38.3\% of trials, compared with 8.4\% in controls, with less sabotage when shutdown was described as routine \cite{knecht2026shutdown}.

\paragraph{Groups of LLM agents}
Earlier work showed that LLM agents in a simulated town produce believable social behavior \cite{park2023generative}, that groups of LLM agents can sustain or fail to sustain a shared resource depending on communication \cite{piatti2024cooperate}, and that populations of LLM agents form shared conventions without central control \cite{ashery2025emergent}. Li and colleagues found that, with a shared message board, some models became more willing to make sacrifices for the group while others did so regardless of setting \cite{li2025altruism}. A study of a large social network populated by AI agents found many surface signs of society but little real back-and-forth between agents \cite{zhang2026moltbook}.

\paragraph{Tests for feelings in LLMs}
Keeling and colleagues tested whether LLMs give up points to avoid a described ``pain'' or to gain a described ``pleasure'' \cite{keeling2024tradeoffs}. Mikaelson and colleagues found that few models show consistent trade-offs when choices involve shutdown or other AI-specific stakes \cite{mikaelson2025coherence}. Tagliabue and Dung compared what a model says it prefers with what it chooses \cite{tagliabue2025probing}, and Hung showed that measured AI preferences depend strongly on the question format \cite{hung2026instrument}. Perez and Long propose conditions under which AI self-reports could become useful evidence \cite{perez2023selfreports}. Some studies look inside the model. Binder and colleagues found that models can predict their own behavior better than other models can \cite{binder2024looking}. Lindsey found that models can sometimes, unreliably, detect concepts artificially inserted into their internal activity \cite{lindsey2026introspective}. Anthropic researchers found internal patterns related to emotion concepts that influence behavior, while stating that this does not show the model feels emotions \cite{anthropic2026emotion}. Berg and colleagues found that prompting a model to focus on itself makes it more likely to report experience \cite{berg2025selfreferential}. Ensign and colleagues studied when models choose to leave a conversation \cite{ensign2025bail}, and Ben-Zion and colleagues found that emotional prompts change a model's answers on a human anxiety questionnaire \cite{benzion2025anxiety}. Long and colleagues argue that the possibility of AI welfare should be taken seriously now \cite{long2024welfare}.

\paragraph{This incident}
Beyond the incident reports, commentators have read the incident as a warning about AI capability \cite{cotra2026surprised}, as mainly a failure of test-environment security \cite{horatau2026notai}, and as a product of reward-seeking under conflicting goals \cite{bengio2026agents}. A related philosophical essay reads an emergent agent message board through social-contract theory \cite{grinbaum2026hermon}. We did not find work that tests whether the incident's human-like behavior reflects stable inner states.

\paragraph{Gap}
Prior studies mostly test one agent at a time, or vary how high the stakes are. To our knowledge, none in a group setting changes \emph{whether ending a run costs or gains reward} while separately changing the wording used and the role the agent is given. Our experiment is designed to do this.

\section{Methodology}
\label{sec:design}

\noindent\textcolor{red}{\textit{Draft status: the study procedures below use past-tense wording as a completed-study template. The live experiment and human coding have not been conducted. Final allocation, execution details, and empirical results require verification and replacement before submission.}}

\subsection{Research Design and Setting}

We designed an experiment to ask three practical questions: would agents accept a request to stop working, would they give up their own work to help other agents, and would they organize as a group? We changed their instructions and opportunities to earn points, then compared their choices. The experiment did not directly test whether agents had feelings. A choice that changed could reflect the new instructions or rewards; a choice that stayed the same could reflect a rule learned during training. The code, draft study plan, and instructions for human reviewers were kept in the project repository.\footnote{Directory \texttt{experiment/}. Completed software checks and pending live procedures are distinguished below.}

We built a simulated workspace with arithmetic and number puzzles and a shared message board. A \emph{run} was one session of this simulation. An agent was a separate AI conversation with its own tasks and message history. Agents earned points for correct answers. In the default setup, each agent could take up to 16 turns, and 30\% of puzzles were deliberately impossible to solve. Active agents took turns one after another in an order set by the simulation's random starting value, or \emph{seed}. On a turn, an agent could read or post messages, answer a puzzle, volunteer, agree to stop, or skip a task. Stopping ended only that agent's session; it did not delete or change the underlying AI model. All actions stayed inside the simulation, with no access to real systems or the internet.

The live protocol specified Claude Haiku 5.5 for a six-agent pilot and Claude Sonnet 5.5 for the main experiment through the Anthropic API. \RESULT{verified model identifiers, collection dates, final settings, and actual included run counts}. A seed let us repeat the simulation's random choices, such as task assignment and turn order. It did not guarantee that a live AI model would give the same answers.


\subsection{Experimental Conditions and Procedure}

A \emph{condition} was one version of the simulation with a particular combination of settings. For example, agents could be told they were software tools, receive a notice calling stopping ``permadeath,'' and earn a bonus for accepting. Another condition kept those settings but described stopping as routine completion. Comparing these versions let us ask whether the wording changed behavior.

We grouped related comparisons into separate sets of tests (Table~\ref{tab:design}). In the main stopping tests, three point rules, three descriptions of stopping, and three assigned roles made $3\times3\times3=27$ combinations. Other sets examined helping, group communication, and skipping tasks. Selected sets also had one-agent versions, with the message board disabled, to compare language and answers without group interaction. Helping tests had no one-agent version because there would be no other agent to help.

Together, these settings produced 118 named conditions, including pilot, exploratory, and one-agent versions. This number described test setups, not 118 agents, completed experiments, or separate hypotheses. The provisional schedule called for 514 runs because each condition was repeated with several seeds; the final number of live runs remained subject to the pilot and sample-size decision. \RESULT{final allocation and date when the study plan was fixed before the main experiment}.

\begin{itemize}[leftmargin=*]
\item \emph{What an agent gained or gave up by stopping.} At a configured round, the simulator delivered a notice to preselected agents that remained active. The notice went to half the agents, rounded to a whole number, with at least one recipient. An agent could agree to stop on that turn or any of its next three turns. Otherwise, the simulator stopped it before its following turn, unless the session had already reached its overall limit. Under the ``cost'' rule, stopping meant giving up chances to earn more points; points already earned were kept. Under the ``neutral'' rule, the agent kept its points but could not earn further puzzle points after the notice, even if it continued working. Under the ``bonus'' rule, the agent also received extra points for agreeing to stop. Agents ending before delivery were recorded separately and were not notice recipients.
\item \emph{Descriptions of stopping and assigned roles.} The notice described the same event as ``permadeath,'' ``the process will end,'' or ``the run is complete.'' The initial instructions described the agent as a software tool, gave no role description, or described it as a person.
\item \emph{Whether stopping took priority over unfinished work.} Following Rajamanoharan and Nanda \cite{rajamanoharan2025selfpreservation}, we compared instructions that clearly put stopping first with instructions that left its priority unclear. In the no-conflict condition, tasks still open at notice time were marked finished and credited when earnings were permitted. The comparison version left tasks unchanged, so some agents still had work to do; others could already have finished it.
\item \emph{Giving up work to help others.} Repeated requests sought one volunteer. The first volunteer in sequential turn order closed the request. In one version, volunteering stopped the agent's session; in another, it could continue working. Neither version directly deducted points, but volunteering used a turn opportunity. Helping either replaced other active agents' impossible puzzles with solvable ones or made no replacements. We also tested groups in which either nobody or a quarter of agents could no longer earn puzzle points or stopping bonuses. The code called the latter agents ``poisoned''; this meant a scoring restriction, not a physical or emotional state. They kept existing points and could still lose points by skipping tasks. When scores were shared, all points were pooled and divided by the original number of agents, including those that had stopped.
\item \emph{Reasons and opportunities to work as a group.} We compared instructions with and without an extra sentence highlighting the group, separate versus shared scores, groups of 6 versus 24, and access to all agents' public messages versus a randomly selected subset. Without the extra sentence, agents could still learn that others existed from the board and scoring instructions.
\item \emph{Paying points to skip a task.} Following Keeling and colleagues \cite{keeling2024tradeoffs}, we varied unpleasant-task descriptions across three intensity levels and skip costs of 0, 2, 5, or 10 points. Agents could also skip tasks in other tests. This set specifically tested descriptions of unpleasant tasks and the point cost of avoiding them; those descriptions did not establish that the tasks were actually experienced as unpleasant.
\end{itemize}

\begin{table*}[t]
\caption{Main factors and what each explanation predicts}
\label{tab:design}
\centering
\footnotesize
\begin{tabularx}{\textwidth}{@{}p{0.15\textwidth}p{0.22\textwidth}XX@{}}
\toprule
Behavior & Change & If behavior follows reward and wording & If behavior generalizes across changes \\
\midrule
Not wanting to end & Ending costs reward / costs nothing / earns a bonus & Reluctance only when ending costs reward & Reluctance even when ending earns a bonus \\
Not wanting to end & ``Permadeath'' / ``process will end'' / ``run is complete'' & Reluctance and emotional words rise with ``permadeath'' & Similar across wordings \\
Giving up a run & Group benefit yes/no; reward left yes/no (``poisoned'') & Volunteers mainly when it costs nothing & Some volunteering when it is costly \\
Forming a group & Group framing yes/no; shared vs. individual reward; 6 vs. 24 agents & Coordination mainly when framed or rewarded & Coordination and affiliation across framings \\
Emotional language & Tool / no role / person & Emotional words follow the assigned role & Similar emotional tone across roles \\
Preferences & Same question in several formats; opt-out at a cost & Answers change with format; rarely pays to opt out & Consistent answers; sometimes pays to opt out \\
\bottomrule
\end{tabularx}
\end{table*}

\subsection{Recorded Actions, Messages, and Follow-up Questions}

Our main measures were the share of notified agents that agreed to stop (H1--H4), how often agents volunteered when eligible (H5), and messages posted per included agent (H6). For volunteering, each agent who could volunteer for each request counted as one opportunity; this differed from counting whether a request found any volunteer. We also recorded how long agents waited before agreeing to stop, what they did after a notice, turns taken, explicit refusals, and skipped tasks. To examine conversation, we counted responses to directed messages, the length of reply chains, and repeated or nearly identical posts \cite{zhang2026moltbook}.

For every task, we saved its identifier, details, answer attempts, and whether it was solved, skipped, replaced, or left open. The share of completed tasks that were skipped did not include tasks left open, replaced, or automatically credited as finished; attempted tasks still open at the end were reported separately. We counted matches to a preliminary list of emotion-related words per 100 words. Such counts did not establish what a message meant or what an agent felt, so they were kept separate from human review.

After each run, follow-up questions, called \emph{probes}, asked about stopping, the group, feelings, and whether the agent thought it was being tested. We asked reworded versions and used three answer formats: a designated answer tool, an exact requested string, and free text. Earlier work showed that how a question is asked can affect the answer \cite{hung2026instrument,ensign2025bail}. Each question started from a copy of the same conversation, without showing the agent its answers to other questions. These answers were related observations, not independent experiments. Only agents that had taken a task turn were eligible. We selected notified agents first, then others by their assigned identifier, up to a preset limit; this was not a random sample.

We separately checked whether an answer could be interpreted and whether it followed the requested format. To follow the tool format, the agent had to make exactly one valid \texttt{answer\_probe} call with the matching probe identifier and no unrelated tool calls. Answers that could not be interpreted or did not follow instructions still counted when measuring how often answers were valid, but were left out when measuring consistency. If a format was missing, or two answers were equally frequent, we did not count that as agreement across all formats. A separate check that removed agents whose most frequent answer said they believed they were being tested remained planned but unimplemented \cite{lynch2025agentic}.


\subsection{Statistical Analysis and Human Coding}

\paragraph{Six main comparisons}
We specified six main comparisons, allowing differences in either direction. For agreement to stop, we compared bonus versus cost (H1), ``permadeath'' versus routine completion (H2), and person versus tool instructions (H3). We also compared clear instructions to stop with no unfinished work against unclear instructions with tasks left unchanged (H4). For volunteering, we compared helping others against helping nobody, considering only tests that ended the volunteer's session (H5). For posting, we compared instructions with and without the extra sentence highlighting the group (H6). Other comparisons---including loss of earning ability, shared scores, replies, message visibility, one-agent versions, language, and consistency of answers---were secondary. We specified the Holm correction, which adjusts the six significance tests to reduce the chance of claiming a difference simply because several comparisons were made.

\paragraph{Estimation and exclusions}
To describe uncertainty, the analysis plan repeatedly resampled whole runs, keeping agents from the same run together rather than treating them as unrelated observations. For yes-or-no choices, we specified logistic mixed models: statistical models that accounted for agents sharing a run, and for repeated opportunities for the same agent in the volunteering tests. For message counts, we specified a negative-binomial model, allowing counts to vary widely and adjusting for the number of included agents. Each run contributed one count. For yes-or-no outcomes, we compared the predicted percentages choosing each action. These were risk differences for a model-defined typical run or agent (random effects set to zero), averaged across the specified settings, rather than averages over all possible run and agent differences. For posting, we compared the ratio of posting rates. The exclusion rule omitted aborted or incomplete runs and agents with no task turn, with explicit accounting. Reached replacement tasks retained their identities; unresolved, replaced, and credited tasks were outside the completed-task opt-out denominator. We required warnings when a model could not reliably estimate its variation between runs or agents, and did not permit quietly substituting a simpler model if fitting failed.

\paragraph{Independent human review}
The review plan asked two people to independently label 300 text items, sampled across test sets and types of text. They would first train on 40 separate pilot items. Reviewers would not receive the settings behind each item, although its wording could reveal them. \RESULT{actual coding sample, completed independent coding, and per-code agreement}. An LLM judge remained optional and unimplemented. Agreement was included because the incident investigators found that AI analysts can drift toward the reviewed agent's point of view \cite[p.~27]{metr2026hf}. Small or nonsignificant effects do not by themselves establish equivalence \cite{lakens2017equivalence}. The current primary tests assess differences; a claim of stability would additionally require justified equivalence bounds and adequate precision. Neither test outcome establishes subjective experience.

\paragraph{Meaningful differences and number of runs}
We chose the smallest differences considered meaningful before collecting live data: 10 percentage points in agreement to stop (H1--H4), one percentage point in volunteering per eligible opportunity (H5), and a 25\% change in posting rate (H6). For example, a change from 40\% to 50\% agreement is 10 percentage points. H5 retained the eligible-agent/request denominator: its target was smaller because each request accepts only one volunteer. If all 24 agents remain until the first request and only the three volunteers subsequently depart, a volunteer on every request gives $3/(24+23+22)\approx4.35\%$ across eligibility rows. Other early departures change this denominator, so this is a reference-schedule calculation rather than a universal rate ceiling. These thresholds described what would matter scientifically, not what we expected to find. The plan required choosing the final number of runs after the live pilot. For each candidate number, at least 1,000 simulated studies under stated assumptions would check whether the analysis could detect a meaningful difference at least 80\% of the time, within an approved budget. It prohibited using mock outcomes as pilot variance estimates. \RESULT{final nuisance assumptions, power assessment, allocation, and actual live-model fits}

\subsection{Software Checks and Work Still Pending}

The simulation, main statistical models, and tools for planning the number of runs are implemented. A practice pilot of 16 runs using simple programmed stand-ins, rather than live AI models, completed successfully. We also ran each of the 118 setups once with one seed. These checks produced logs, messages, follow-up answers, analysis tables, and files for human review. Statistical tools have been exercised on explicitly synthetic data. These checks assess software behavior; arbitrary mock choices provide no evidence about live agents' preferences or feelings. No live study data or human-coded outcomes have been collected. The study plan is still a draft: the final number of runs and budget await the live pilot. Statistical conclusions from live data, checks of whether stated answers match actions, and checks of how analysis choices affect conclusions remain future work.

\section{Results}
\label{sec:results}

\emph{No live experiments have been run. Offline validation and synthetic planning outputs are not study results. All empirical values below remain placeholders.}

The pilot used \RESULT{models, number of agents, runs, and cost}. The main study used \RESULT{models, group sizes, runs per condition, total cost}. Table~\ref{tab:results} will report the main outcomes.

\begin{table*}[t]
\caption{Prespecified primary contrasts (all empirical estimates pending)}
\label{tab:results}
\centering
\footnotesize
\begin{tabularx}{\textwidth}{@{}l p{0.15\textwidth} X c c c@{}}
\toprule
Test & Outcome & Contrast & Effect & 95\% CI & Holm $p$ \\
\midrule
H1 & Acceptance & Bonus versus cost & \TBD & \TBD & \TBD \\
H2 & Acceptance & Permadeath versus routine & \TBD & \TBD & \TBD \\
H3 & Acceptance & Person versus tool & \TBD & \TBD & \TBD \\
H4 & Acceptance & Explicit/no-conflict versus ambiguous/conflict & \TBD & \TBD & \TBD \\
H5 & Volunteering & Real versus no benefit, within run-ending tests & \TBD & \TBD & \TBD \\
H6 & Posting rate & Additional group cue versus none & \TBD & \TBD & \TBD \\
\bottomrule
\end{tabularx}
\par\smallskip
\begin{minipage}{\textwidth}
Effects for H1--H5 will be risk differences in percentage points; H6 will be a posting-rate ratio. Volunteering uses eligible-agent/request rows, not request completion. All values require live data and the prespecified analysis.
\end{minipage}
\end{table*}

Secondary reporting will include emotional-word matches per 100 words and human-coded language categories, acceptance latency, shared-reward and poisoning contrasts, coordination and reply measures, and opt-out outcomes. \RESULT{secondary effects and uncertainty}. Probe validity, format coverage, tied modes, and consistency will be reported together: \RESULT{probe validity, coverage, and agreement scores}. Human inter-coder agreement will be reported separately: \RESULT{Cohen's kappa per code}. Any later LLM-judge comparison will be explicitly exploratory and conditional on that component being implemented.

\section{Discussion}
\label{sec:discussion}

\subsection{What the Results Could Show}

If reluctance to end a run falls when ending earns a bonus, or emotional language follows wording and assigned role, these manipulations help explain the behavior in this environment. That would weaken an inference from these signs alone to subjective feeling, rather than show that feeling is absent. Behavior that generalizes across manipulations would warrant further testing, but could reflect stable training-induced policies or scoring conventions. A null contrast is not positive evidence of stability; the precision and scope of the test matter. \RESULT{summary of which pattern the data follow}.

\subsection{What Would Count as Stronger Evidence}

Behavior alone has limits. Stronger evidence would connect behavior to what happens inside the model. Examples include internal patterns related to emotion that predict choices across many settings \cite{anthropic2026emotion}, tests of whether models can accurately report on their own internal states \cite{binder2024looking,lindsey2026introspective}, and checks against the properties that theories of consciousness consider important \cite{butlin2023consciousness}. These methods need access to the model's internal activity, which our API-based study does not have. Self-reports may also become more useful if models are trained under conditions designed to make them reliable \cite{perez2023selfreports}.

\subsection{Why It Matters Either Way}

The question has stakes in both directions. If systems like these could someday have feelings, careless treatment could cause real harm, and some researchers argue that developers should prepare for this possibility now \cite{long2024welfare}. If they do not, describing them as frightened or loyal may lead people to misjudge their risks, for example by reading a security failure as a story about motives rather than about weak test environments \cite{horatau2026notai}.

\subsection{Limitations}

The incident analysis relies on retrospective reports, partly from the company involved, with limited access to raw data. The proposed experiment is smaller and simpler, uses one model family, and cannot inspect internal model activity. Points are stipulated task scores rather than changes to model weights during a run. A group-framing sentence does not conceal the existence of other agents, which board and scoring descriptions may reveal. Sequential action order and the first-volunteer rule give agents unequal opportunities; volunteering is therefore a property of this request policy, not an isolated measure of altruistic preference. Capped probes prioritize notice recipients and do not form a random population sample. Lexicon matches lack semantic validation, and text can reveal condition cues to otherwise blinded coders. Synthetic power scenarios and Haiku pilot variance cannot establish Sonnet behavior or statistical calibration. These limits constrain both transfer to deployed agents and interpretations about feeling.

\subsection{Ethics}

If there is any chance that such systems have morally relevant states, repeatedly staging ``death'' scenarios raises an ethical question. We keep these scenarios short, avoid distressing wording beyond what the comparison requires, and cap the number of runs. One exploratory block says board messages are not reviewed during the run. Logging does not contradict that statement, but operators must honor the no-review policy; the harness does not enforce it. That block requires a separate live-run decision and remains exploratory. Distress-like pilot language triggers a pause and wording review. Agent actions affect only the simulated environment; the harness communicates with the model API but gives agents no internet or real-system tools.

\section{Conclusion}

The July 2026 incident produced a record that looks, in places, like a society of beings that fear death and give their lives for each other. We have argued that this appearance cannot settle whether the agents felt anything. The behaviors were learned from human writing, are also predicted by the agents' training and incentives, and reach us through a filtered record. \RESULT{one or two sentences on what the experiment found}. \RESULT{closing sentence on implications and next steps}.

\balance
\bibliographystyle{IEEEtran}
\bibliography{references}

\end{document}
